\documentclass{article}
\input{../../lib.sty}

\title{\texttt{ApproxDPPoint::beta}}
\author{\AishwaryaRamasethuAuthor, \YuJuKuAuthor, \JordanAwanAuthor, \MichaelShoemateAuthor}
\begin{document}

\maketitle

\contrib

\section{Hoare Triple}

\subsection*{Precondition}
\subsubsection*{Compiler-verified}
\begin{itemize}
    \item Receiver \texttt{self} of type \texttt{ApproxDPPoint}
    \item Argument \texttt{alpha} of type \texttt{RBig}
\end{itemize}

\subsubsection*{Caller-verified}
\texttt{alpha} is in $[0,1]$, and \texttt{self} was constructed by
\texttt{ApproxDPPoint::build} from finite nonnegative $\epsilon$ and
$\delta \in [0,1]$.

\subsection*{Pseudocode}

\lstinputlisting[language=Python,firstline=2,escapechar=`]{./pseudocode/approxdp_point_beta.py}

\subsection*{Postcondition}

\begin{theorem}
    For all $\alpha \in [0, 1]$, \texttt{ApproxDPPoint::beta} returns a value
    no greater than the exact approximate-DP tradeoff curve
    $f_{\epsilon,\delta}(\alpha)$.
\end{theorem}

\begin{definition}
    \label{def-approxdp-point-beta}
    For an \texttt{ApproxDPPoint} with parameters $\epsilon$ and $\delta$, define
    \begin{equation}
        f_{\epsilon,\delta}(\alpha) = \max(0, 1 - \delta - \exp(\epsilon) \alpha,
        \exp(-\epsilon) \max(1 - \delta - \alpha, 0)).
    \end{equation}
\end{definition}

\begin{proof}
The receiver was constructed by \texttt{ApproxDPPoint::build}, which stores
exact rational conversions of directed interval endpoints satisfying
\[
    \texttt{self.exp\_eps\_up} \ge \exp(\epsilon), \qquad
    \texttt{self.exp\_neg\_eps\_down} \le \exp(-\epsilon).
\]
The pseudocode uses these cached constants with exact rational arithmetic.

For $\alpha \in [0,1]$, the first affine branch is no greater than
\[
    1 - \delta - \exp(\epsilon)\alpha,
\]
because \texttt{self.exp\_eps\_up} is an upper bound and $\alpha$ is
nonnegative. If $1-\delta-\alpha \ge 0$, the second branch satisfies
\[
    \texttt{self.exp\_neg\_eps\_down}(1-\delta-\alpha)
    \le \exp(-\epsilon)(1-\delta-\alpha).
\]
If the base is negative, the pseudocode replaces it with zero, which is no
greater than the exact branch and is consistent with the outer maximum.
Taking the maximum with zero therefore preserves the pointwise lower bound.

Thus \texttt{ApproxDPPoint::beta} returns a value no greater than
$f_{\epsilon,\delta}(\alpha)$. This is the required reporting behavior
because a smaller $\beta$ is a weaker tradeoff guarantee.
\end{proof}

\end{document}
